\documentclass[11pt]{article}

\usepackage[T1]{fontenc}
\usepackage[utf8]{inputenc}
\usepackage{lmodern}
\usepackage[margin=1in]{geometry}
\usepackage{microtype}
\usepackage{enumitem}
\usepackage{tabularx}
\usepackage{pdflscape}
\usepackage[round,authoryear]{natbib}
\usepackage[hidelinks]{hyperref}

\title{Axodus: A Prototype-Stage Conceptual Architecture for Traceability and Accountability in Federated Modular Digital Institutions}
\author{Mauricio Z. Filho\\
\small Independent Researcher, Axodus\\
\small Correspondence: mzfshark@gmail.com\\
\small Project discussion: Telegram @thinkincoin\\
\small Author handle: @mzfshark}
\date{Draft v0.6}

\begin{document}
\hypersetup{
  pdftitle={Axodus Architecture Paper v0.6 - English Controlled Review Copy},
  pdfauthor={Mauricio Z. Filho},
  pdfsubject={CC BY 4.0 manuscript prepared for a private controlled academic-review process},
  pdfkeywords={conceptual architecture, controlled academic review, English authoritative manuscript}
}
\maketitle

\begin{center}
\fbox{\parbox{0.88\linewidth}{\centering
\textbf{Status:} Internal v0.6 academic draft prepared as a private controlled human-review copy; reviewer delivery is not authorized, and external circulation and publication remain blocked.\\
\textbf{Document type:} Prototype-stage conceptual architecture and focused research-agenda paper.\\
\textbf{License:} CC BY 4.0.}}
\end{center}

\begin{center}
\fbox{\parbox{0.92\linewidth}{\centering
\textbf{CONTROLLED REVIEW COPY - ENGLISH AUTHORITATIVE MANUSCRIPT}\\
This manuscript is provided through a private, controlled academic-review
process and is licensed under CC BY 4.0. Package ID:
\texttt{AXODUS-V0.6-CHR-EN-20260712}. Reviewers are requested, as a matter of
professional courtesy and process integrity, not to redistribute or publicly
discuss the manuscript before the review process closes. This procedural
request does not modify or restrict the rights granted under CC BY 4.0.
Reviewer delivery requires separate authorization. Controlled review does not
imply endorsement, validation, acceptance, or peer-review completion.}}
\end{center}

\begin{abstract}
Human-AI-assisted digital institutions become difficult to govern when
constitutional constraints, bounded local governance domains, modular
capabilities, and review obligations must be coordinated within the same
decision environment. Adjacent literatures on socio-technical systems,
decentralized governance, provenance, accountable algorithms, contestability,
and human-AI interaction provide differently scoped components for examining
this problem. This paper presents Axodus as one proposed integration of those
components: a prototype-stage conceptual architecture for that setting. Its
central contribution is a six-layer governance-oriented model relating a
constitutional center, bounded local governance domains, modular
institutional capabilities, evidence and provenance, bounded AI assistance,
and risk or contestability controls. The six layers are presented as one
defensible conceptual decomposition, not as a necessary, minimal, uniquely
correct, or universally sufficient structure. Traceability and
reconstructability may support accountability and contestability, but do not
establish either property independently. Axodus is described only as a
conceptual, prototype-stage model; the paper reports no deployment, empirical
validation, observed institutional use, or operational readiness.
\end{abstract}

\noindent\textbf{Keywords:} socio-technical systems; federated digital
institutions; human-AI coordination; traceability; accountability;
governance; contestability

\section{Introduction}

Digital institutions increasingly combine distributed authority, modular
capabilities, public records, and automated assistance. Governance becomes
harder when these elements must remain reviewable across a constitutional
center and bounded local governance domains rather than inside a single
authority chain. In that setting, the institutional challenge is not only
whether a decision can be documented, but whether later reviewers can
reconstruct which authority applied, which bounded capability was invoked,
what evidence was considered, how AI assistance shaped the record, and how
correction or challenge could proceed. Socio-technical systems
research treats those questions as jointly organizational and technical rather
than reducible to one domain alone \citep{baxter2011sociotechnical}.

Relevant literatures address differently scoped and partially overlapping
parts of this problem. Research on decentralized governance shows
that self-executing rules do not remove interpretation, amendment, or social
coordination \citep{hassan2021dao,hsieh2018bitcoin}. Human-AI interaction
guidance emphasizes review, correction, and expectation setting rather than an
undifferentiated transfer of authority \citep{amershi2019guidelines}.
Provenance and credentialing standards help distinguish evidence lineage and
machine verification from institutional judgment
\citep{w3c2013provdm,w3c2025vcdm}. This paper does not claim an exhaustive
demonstration that no comparable architecture exists. It develops a bounded
integrative synthesis for a specific institutional problem.

This paper introduces Axodus as a prototype-stage conceptual architecture for
human-AI-assisted digital institutions under bounded governance constraints.
The project originates from a multi-year independent research and design
effort; that statement records background provenance rather than scientific
evidence. The present contribution remains conceptual: it does not report a
platform in real institutional use, a functioning institutional environment,
evaluation results, or independently tested controls.

The contribution is architecture-level design reasoning rather than an
implementation claim. The paper proposes a layered model for how a
constitutional center, bounded local governance domains, modular
institutional capabilities, evidence packages, attributable AI assistance,
and review paths may be arranged into one inspectable decision structure. The
six-layer structure is one defensible conceptual decomposition among possible
alternatives. Participation readiness remains a bounded supporting concern,
while distinctions among documentation, development, production, and
authority serve only as a scope boundary.

Implementation, deployment, institutional use, and empirical validation are
outside the paper's scope. The manuscript also excludes token design,
financial operations, educational-distribution mechanics, security
guarantees, and legal conclusions.

\section{Background and Related Work}

\subsection{Socio-technical systems and federated digital institutions}

A socio-technical perspective asks how roles, incentives, work practices,
information, and technical components jointly shape outcomes. Baxter and
Sommerville argue for integrating organizational concerns with systems
engineering rather than treating them as a late adoption problem
\citep{baxter2011sociotechnical}. Axodus adopts this perspective by modeling
governance, accountability, and review as architectural elements rather than
external policies attached after implementation.

The term \emph{federated modular digital institution} is used here for an
organized system in which a central constitutional layer defines shared
constraints, while bounded local governance domains may configure modular
institutional capabilities. This use is conceptual. It does not imply a
deployed federation, legal federal status, or operating multi-tenant
platform.

\subsection{Decentralized governance, modular coupling, and review}

Hassan and De Filippi describe a DAO as a blockchain-based system in which
people coordinate through decentralized governance and self-executing rules
\citep{hassan2021dao}. Hsieh et al.\ distinguish machine consensus from the
social consensus required to alter protocols and organizational arrangements
\citep{hsieh2018bitcoin}. These distinctions caution against equating
technical distribution with accountable governance.

Smart-contract-based financial markets provide one domain-specific example in
which composability can propagate dependencies across connected protocols
\citep{schar2021defi}. Systems-theoretic safety research instead offers the
broader claim needed here: failures in complex socio-technical systems can
arise from interactions, control mismatches, and poorly governed coupling
rather than from isolated component defects alone
\citep{leveson2012engineering}. The present paper therefore treats that
domain only as one motivating example of dependency propagation, not as the
paper's institutional center.

\subsection{Human-AI oversight and accountable decision support}

Guidelines for human--AI interaction emphasize communicating system
capabilities, supporting correction, and learning from behavior over time
\citep{amershi2019guidelines}. At the organizational level, the NIST AI Risk
Management Framework treats governance, mapping, measurement, and management
as connected risk functions \citep{tabassi2023airmf}. These sources motivate
an architecture in which AI outputs remain attributable, reviewable inputs to
human decisions rather than sovereign decisions in themselves.

Constitutional AI uses written principles to guide model behavior through
supervised learning and AI feedback \citep{bai2022constitutional}. Axodus uses
\emph{constitutional} differently: it refers to institutional rules about
authority, prohibitions, amendments, review, and challenge. The proposed model
does not claim to train or align an AI model through Constitutional AI.

\subsection{Provenance, credentials, and institutional judgment}

The W3C PROV Data Model describes provenance through entities, activities,
agents, derivations, attribution, and responsibility relationships
\citep{w3c2013provdm}. This vocabulary is useful for distinguishing source
artifacts, transformations, automated assistance, and responsible roles in a
decision record. It supports reconstruction of how an output was produced,
but does not by itself establish that the decision was legitimate,
contestable, or accountable.

The W3C Verifiable Credentials Data Model separates the cryptographic
verification of a credential from validation of whether its claims are fit for
a verifier's purpose \citep{w3c2025vcdm}. This matters here because machine-
verifiable eligibility artifacts would not by themselves establish fair
assessment, legitimate authorization, or accountable participation
boundaries.

Work on accountable algorithms likewise argues that automated decision
processes must be assessable against external legal or policy standards and
that transparency alone is insufficient \citep{kroll2017accountable}.
Provenance is therefore one input to accountability, not its substitute.

\subsection{Contestability and reviewable governance}

Contestability-by-design literature identifies safeguards, human review and
intervention, and third-party scrutiny as design features for challenge and
correction \citep{alfrink2023contestable}. This framing is directly relevant
to federated modular governance because a traceable record remains
insufficient if a person affected by a decision cannot identify the
applicable rule, challenge evidence, or reach an appropriately independent
reviewer.

\subsection{Conceptual synthesis and architecture derivation}

The architecture is derived through a narrative, non-systematic conceptual
synthesis rather than a systematic review, formal proof, or empirical design
method. The synthesis treats the literature strands in Sections 2.1--2.5 as
conceptual foundations, the five tensions in Section 3 as design problems, and
the resulting requirements as reasons for separating responsibilities across
six layers. Layer boundaries are drawn where authority, local interpretation,
bounded capability, record reconstruction, advisory analysis, and challenge
conditions require different responsibilities or expressly excluded powers.

This procedure does not establish that six layers are necessary, minimal,
uniquely correct, or universally sufficient. It presents one defensible
conceptual decomposition whose boundaries can be inspected, compared with
alternatives, and revised or rejected. The derivation matrix makes the chain
from foundations to evaluation constructs explicit; it does not imply that
the cited literature validates Axodus.

\begin{landscape}
\begin{center}
\scriptsize
\begin{tabularx}{\linewidth}{|X|X|X|X|X|X|}
\hline
\textbf{Foundation} & \textbf{Tension} & \textbf{Requirement} &
\textbf{Layer(s)} & \textbf{Property supported} &
\textbf{Evaluation construct} \\
\hline
Socio-technical systems and institutional governance & Central constraints
versus local autonomy & Separate higher-order constraints from bounded local
interpretation & Constitutional governance; bounded local governance domain &
Inspectable authority and answerability conditions & Authority identification;
conflict and scope detection \\
\hline
Modularity and systems coupling & Modular capabilities versus governability &
Declare capability boundaries, dependencies, and shared design constraints &
Functional service module; constitutional governance & Cross-layer coherence
and bounded responsibility & Dependency completeness; unsupported-transition
detection \\
\hline
Provenance and accountable decision support & Documentation versus
institutional reality & Attribute evidence, transformations, roles, and
dispositions without equating records with legitimacy & Evidence and
provenance & Traceability, reconstructability, and reviewability &
Reconstruction completeness; provenance visibility; missing-evidence detection \\
\hline
Human-AI interaction and AI risk governance & Automation versus accountable
disposition & Keep AI outputs attributable and advisory, with explicit human
disposition & AI assistance; evidence and provenance & Bounded reliance and
attributable responsibility & Error detection; override behavior; rationale
quality \\
\hline
Contestability-by-design & Authority or evidence disputes versus effective
review & Preserve challenge triggers, independence conditions, and reachable
review paths & Risk boundary and contestability; constitutional governance &
Contestability-supporting and corrective conditions & Challenge-path
accessibility; reviewer disagreement; disposition classification \\
\hline
Participation and readiness boundary & Participation versus bounded readiness &
Treat readiness as a contestable, accommodated input rather than universal rank
or authority & Bounded local governance domain; risk boundary and
contestability & Reviewable participation conditions & Exclusion-risk analysis;
accommodation and challenge visibility \\
\hline
\end{tabularx}
\end{center}
\end{landscape}

\section{Problem Statement}

The architecture addresses a design problem rather than an established
empirical effect:

\begin{quote}
How can a constitutionally constrained federated modular digital institution
support traceability, reconstructability, and the institutional conditions
for accountability and contestability when bounded local governance domains
configure modular institutional capabilities through human-AI-assisted
decision processes?
\end{quote}

Five tensions structure the problem.

\begin{enumerate}[leftmargin=*]
\item \textbf{Central constraints versus local autonomy.} A constitutional
layer may preserve coherence, while bounded local governance domains still
require meaningful discretion over local decisions and configurations.
\item \textbf{Modular capabilities versus governability.} Service separation
may limit direct coupling, while modular selection can create hidden
dependency chains, ambiguous authority, and fragmented review paths.
\item \textbf{Automation versus accountable disposition.} AI assistance can
reduce analytical load, but responsibility can become unclear when its
outputs influence authorization, escalation, or review.
\item \textbf{Participation versus bounded readiness.} Open participation may
support legitimacy, while consequential governance actions may require
reviewable preparation, accommodations, and challenge.
\item \textbf{Documentation versus institutional reality.} A detailed
specification can improve inspection while also creating a misleading
impression of implementation completeness or legitimate authority.
\end{enumerate}

Axodus is proposed as one defensible conceptual decomposition for making these
tensions inspectable and reviewable. It is not presented as a complete,
optimal, uniquely correct, or empirically tested solution.

\section{Conceptual Architecture}

The model comprises six logical layers selected through the derivation in
Section 2.6. They are conceptual design surfaces for allocating authority,
structuring records, and preserving reviewability, not an implementation
topology. Their separation expresses distinct responsibilities and excluded
powers; it does not assert that every institution requires this decomposition.

Information and control move through the architecture in different ways.
Control begins with constitutional constraints that define what local
decisions, capability selections, and escalation paths are in scope.
Evidence packages, provenance records, advisory analysis, rationale, and
review records then make those decisions reconstructable for later review.
Risk boundaries constrain permissible state transitions, while contestability
paths preserve challenge, correction, escalation, and review conditions.

\noindent\textbf{Illustrative synthetic scenario - bounded-governance
exception review.} A mock bounded local governance domain requests a
temporary exception affecting one modular institutional capability. The
request includes the applicable constitutional rule, local authority basis,
declared conflicts, and a controlled evidence package. A simulated AI review
flags one conflict but omits a dependency between the requested
configuration and a review requirement. The designated human reviewer
withholds disposition, records the omission, and escalates the matter because
a risk boundary is triggered. An independent review role can then reconstruct
which constitutional constraint applied, what local choice was proposed, what
the simulated AI review changed, and why escalation occurred. This scenario
is an analytical illustration, not evidence of an Axodus implementation or
institutional process.

\begin{enumerate}[leftmargin=*]
\item constitutional governance defines the applicable roles, decision
rights, prohibitions, amendment rules, and challenge conditions;
\item bounded local governance domains interpret that scope for a local
decision context and may propose a configuration or exception within
constitutional limits;
\item modular institutional capabilities express what is being selected,
reviewed, or changed;
\item the evidence and provenance layer links source materials, authority
statements, conflicts, transformations, and responsible roles;
\item human-AI assistance may analyze that material, but only as attributable
advisory input to human disposition; and
\item risk-boundary and contestability controls keep challenge, correction,
escalation, and review conditions available.
\end{enumerate}

The layer-interface table specifies conceptual responsibilities and handoffs.
Its inputs and outputs are record categories, not APIs, software components,
or claims about implemented infrastructure.

\begin{landscape}
\begin{center}
\scriptsize
\begin{tabularx}{\linewidth}{|X|X|X|X|X|X|X|}
\hline
\textbf{Layer} & \textbf{Responsibility} & \textbf{Input} & \textbf{Output} &
\textbf{Dependency} & \textbf{Boundary} & \textbf{Excluded authority} \\
\hline
Constitutional governance & Define higher-order roles, constraints, amendment,
escalation, and challenge conditions & Institutional rules and amendment basis &
Applicable constraint and authority record & Constrains local domains, modules,
and review conditions & Non-delegable constraints cannot be expanded locally;
rules remain inspectable & No claim of legal force, deployment, infallibility,
or automatic enforcement \\
\hline
Bounded local governance domain & Interpret permitted scope for a local
decision context & Constitutional constraints and local proposal & Scoped
decision, configuration, or escalation request & Depends on constitutional
scope and module declarations & Cannot override non-delegable constraints & No
sovereign, unlimited, or production authority \\
\hline
Functional service module & Express a bounded institutional capability and its
dependencies & Authorized scope, task, evidence requirements, and dependency
declarations & Capability-specific proposal, result, or failure record &
Constrained by governance; inspected through evidence and risk layers & Must
expose purpose, dependencies, evidence needs, escalation, and failure
conditions & Does not own constitutional authority, provenance, AI disposition,
or contestability \\
\hline
Evidence and provenance & Attribute sources, transformations, agents,
assistance, and dispositions & Source artifacts, authority statements,
analyses, conflicts, and decisions & Linked decision-and-evidence record &
Receives records from all decision layers and supports later review & Record
linkage does not establish truth, integrity, legitimacy, or accountability &
No automatic validation, judgment, or integrity guarantee \\
\hline
AI assistance & Provide bounded, attributable advisory analysis & Declared
task, source material, and constraints & AI output with provenance and human
disposition & Depends on evidence inputs and human review; subject to risk
boundaries & Human authority and responsibility remain explicit & No sovereign
decision, enforcement, or binding judgment \\
\hline
Risk boundary and contestability & Surface review triggers, challenge
conditions, independence needs, and possible dispositions & Proposed action,
applicable boundary, disputed evidence, and conflicts & Review, refusal,
escalation, correction, reversal, or no-change record & Depends on
constitutional scope and reconstructable evidence & A challenge path must be
identifiable and separable from the original disposition & No claim of standing
procedure, effective remedy, or guaranteed accountability \\
\hline
\end{tabularx}
\end{center}
\end{landscape}

\subsection{Constitutional governance layer}

The constitutional governance layer is a conceptual constraint surface that
defines higher-order rules: roles, decision rights, prohibitions, amendment
procedures, escalation paths, and challenge conditions. Its purpose is to make
non-delegable institutional constraints visible before local configuration
begins and to frame how later review can distinguish a permitted action from
an out-of-scope one.

In this proposal, the layer structures accountability-supporting boundaries rather than
serving as a deployed constitution, legal instrument, active authority, or
enforcement body. A reviewable action therefore depends on both institutional
justification and whatever technical authorization a future implementation
might use; neither alone would establish legitimacy.

Because constitutional constraints can themselves be incomplete, ambiguous, or
overbroad, they must remain inspectable design objects with bounded review,
amendment, and challenge conditions. The model therefore does not treat the
constitutional layer as infallible or beyond contestation.

\subsection{Bounded local governance domain layer}

As a conceptual reference rather than an implementation equivalence,
polycentric governance describes multiple decision centers that may be
formally distinct yet interdependent and may operate across nested governance
layers \citep{ostrom2010beyond}.
The bounded local governance domain layer represents a locally scoped
decision context operating under constitutional constraints. A bounded local
governance domain may configure local processes, select functional service
modules, and allocate decision tasks only within the scope permitted by the
constitutional layer. Axodus institutional documentation may use the term
Tenant for such bounded local governance domains; this paper uses the
governance-domain phrasing to avoid implying deployed SaaS tenancy, active
customer entities, operating communities, or production infrastructure.

A bounded local governance domain cannot expand or override non-delegable
constitutional constraints. When the local interpretation of a proposed
action conflicts with a constitutional boundary, the architecture treats that
ambiguity as a review condition requiring documentation and, where
necessary, escalation before final disposition.

\subsection{Functional service module layer}

As a conceptual reference rather than validation of this taxonomy, modularity
literature describes complex systems built from smaller subsystems that can be
designed separately while functioning together under shared design rules
\citep{baldwin2000design}.
Functional service modules are abstract bounded institutional capabilities
through which selected functions can be configured, limited, and governed.
The paper does not present a product catalog. Instead, it treats a
functional service module as a conceptual unit that should declare its
purpose, inputs, outputs, authorized scope and decision dependencies, evidence requirements,
dependencies, escalation conditions, and failure behavior.

A functional service module is therefore a bounded capability surface, not
the owner of constitutional authority, evidence provenance, AI assistance,
or contestability. Those remain separate cross-cutting layers that
constrain, inspect, and reconstruct module-specific decisions.

Axodus institutional documentation may use the term service nuclei in broader
institutional contexts; this paper uses functional service modules as the
architecture-facing abstraction to avoid implying a product catalog,
implementation inventory, or deployed system decomposition.

A minimal functional taxonomy may include authority and role registries,
evidence registries, decision-rationale records, conflict and disclosure
registers, review and challenge paths, risk-boundary controllers, and
participation-readiness support as component types associated with
module-governed capabilities. These are component types rather than claims
about a definitive Axodus module inventory or implementation state.

\subsection{Evidence and provenance layer}

The evidence and provenance layer links proposals, supporting materials,
authority statements, declared conflicts, AI-generated analyses, decisions,
and subsequent reviews. Its design objective is reconstructability: a
reviewer should be able to ask what was known, who was authorized, what
configuration was proposed, what assistance was used, and why an outcome
followed.

At the conceptual level, source documents and outputs can be treated as
entities, analysis and review steps as activities, and human or automated
participants as agents. Derivation and attribution relations can then connect
an output to its inputs and responsible roles \citep{w3c2013provdm}. These
concepts specify what a reviewable record should express so that a later
reviewer can reconstruct the path from source material to disposition. They
do not by themselves prove that a record is true, complete, tamper-resistant,
or accountable in the stronger normative sense.

Accordingly, the layer supports later inspection rather than automatic
validation. It does not prescribe a storage system, cryptographic proof,
implemented evidence infrastructure, or integrity guarantee, and it does not
convert record linkage into a full judgment about whether the decision was
substantively justified.

\subsection{AI-assistance layer}

AI assistance remains advisory. It may support bounded tasks such as document
comparison, proposal summarization, inconsistency detection, constitutional
analysis support, adversarial checking, and planning-oriented review support.
Each use should declare the task, source material, output, uncertainty where
available, human reviewer, and disposition so that later reviewers can
inspect how advisory analysis entered the record. Humans retain decision
authority and responsibility.

This layer is intentionally role-abstract. The paper does not use named
agent labels, does not claim autonomous execution authority, and does not
claim that AI assistance itself enforces governance, issues binding
judgments, or replaces human institutional review.

\subsection{Risk-boundary and contestability layer}

The risk-boundary and contestability layer is a conceptual surface for making
it visible when further review, refusal, documentation, or a challenge path
may be needed. In this model, a boundary should identify the relevant role,
the review trigger, the evidence required for inspection, and the condition
under which a proposed action exceeds its permitted scope. Contestability
therefore concerns whether a decision can be reconstructed, questioned, and
re-examined in principle, not whether the architecture already operates a
standing challenge mechanism.

A minimally challengeable record should identify the contested decision or
record, the applicable rule or boundary, the disputed evidence or rationale,
the review-independence condition expected for re-examination, and the
possible disposition category, such as correction, escalation condition,
refusal, reversal, or no change. These are conceptual fields for
reconstructability and challengeability, not a standing challenge procedure.

Within that framing, reviewer independence is a design concern rather than a
claim about a standing body or named role. The architecture should make it
possible to distinguish the original decision context from later challenge or
re-examination conditions, including disclosed conflicts and situations in
which the same actor should not be the sole judge of a contested record.
Human oversight becomes nominal when a reviewer lacks the time, authority
context, evidence access, or practical capacity to change an outcome
\citep{tabassi2023airmf,amershi2019guidelines}. Contestability-by-design
literature is relevant here because it treats review visibility,
intervention, and external scrutiny as design requirements for challengeable
systems rather than as guarantees that harmful outcomes are prevented
\citep{alfrink2023contestable}.

\subsection{Supporting participation-readiness boundary}

Participation readiness is retained only as a bounded supporting concern for
preparation, accommodation, review, and challenge in selected consequential
contexts. It is not a co-equal contribution, a general social ranking system,
or a universal competence filter. Any readiness artifact would be a
contestable and reviewable input, not an independent source of legitimacy,
fairness, or authority. Criteria would need to remain proportionate,
accommodated, and open to challenge; exclusion, surveillance, interpretive
capture, test gaming, and privacy risk remain reasons to narrow, redesign, or
reject the mechanism. Its validity is future research rather than a claim
evaluated in this paper.

\section{Prototype Status}

Axodus is presented in this paper as a conceptual, prototype-stage model. The
public repository presently supports editorial artifacts, governance records,
and research documentation. It does not supply evidence of an integrated
implementation, functioning federated platform, production operation, user
adoption, empirical performance, or effective controls.

The absence of those claims is substantive. Architectural consistency can be
reviewed before implementation, but it cannot substitute for observed
behavior. Future versions should maintain a component-by-component evidence
table distinguishing designed, prototyped, tested, deployed, and
independently evaluated states.

Documentation maturity, development maturity, production readiness, and
institutional authority are separate determinations. This distinction is used
only to prevent category error: documentation does not establish
implementation, implementation does not establish production authorization,
and maturity does not establish legitimate authority or operational
capability. The paper does not retain L-Level or D-Level as standalone
scientific constructs, assign levels to components, or present a maturity
taxonomy as a validated instrument.

Scope reduction remains a design requirement. Initial prototypes should focus
on decisions for which reconstruction and challenge are materially important,
use a bounded decision-and-evidence record, separate advisory AI outputs from
human dispositions, and remove elements that do not improve inspectability.

\section{Evaluation Plan}

The evaluation program is organized around the integrated architecture while
remaining controlled, conceptual, synthetic, and artifact-based. It evaluates
whether the proposed decomposition makes authority, dependencies, records,
advisory assistance, and challenge conditions more inspectable. It does not
evaluate an operating institution or establish legitimacy, fairness,
accountability, contestability effectiveness, or remedy.

The initial study stage would use synthetic bounded-local-governance
scenarios, controlled evidence packages, simulated AI-assistance records,
planted errors, and a high-fidelity mock-up or wizard-of-oz review interface
rather than an active institutional setting. A preregistered protocol would
define reviewer tasks and study criteria. These are proposed designs, not
reported results.

\paragraph{RQ1: Traceability and reconstruction.} Does the proposed layered
decision-and-evidence structure help independent reviewers reconstruct bounded
local governance decisions involving functional service modules more
accurately and efficiently than a document-only baseline? Illustrative
criteria include reconstruction accuracy and completeness, provenance
visibility, authority-conflict detection, missing-evidence detection, time to
reconstruct rationale, and reviewer disagreement.

\paragraph{RQ2: Cross-layer coherence and review conditions.} Can independent
reviewers use the six-layer representation to identify missing or conflicting
authority, undeclared dependencies, unsupported transitions between layers,
and whether an accessible challenge path and review-independence condition are
specified? Illustrative criteria include cross-layer consistency, dependency
completeness, unsupported-transition detection, challenge-path accessibility,
and correct classification of possible disposition categories. These measures
test inspectability and accountability- or contestability-supporting
conditions; they do not demonstrate institutional answerability, effective
challenge, correction, reversal, or remedy.

\paragraph{RQ3: Human-AI reliance in advisory review.} For bounded
proposal-analysis tasks, how do provenance and mandatory human disposition
affect error detection, calibrated reliance, override behavior, rationale
quality, and review quality? Simulated AI records would include planted errors,
ambiguous cases, and contested recommendations. Illustrative criteria include
planted-error and unsupported-output detection, override or acceptance
behavior by condition, rationale quality, reviewer disagreement, and
correction after provenance information is shown.

For this paper, a document-only baseline uses static policy and record
documentation without the proposed layered review architecture, explicit
provenance structure, or bounded challenge-path framing. Negative findings
matter as much as positive ones: the representation may fail to improve
reconstruction, may increase workload without improving error detection, may
leave cross-layer inconsistencies undetected, may produce disagreement about
challenge conditions, or may not improve calibrated reliance.

\paragraph{Staged future evaluation.} Evaluation should progress from
controlled artifact inspection to bounded reviewer studies using synthetic
scenarios, then to controlled prototype evaluation under a separately approved
protocol. Any later institutionally situated study would require additional
governance, ethics, legal, privacy, and evidence gates. Participation-readiness
validity and broader institutional accountability remain later research
topics.

\section{Limitations and Threats to Validity}

This paper is conceptual and still broader than a mature external article
would ideally be. The architecture can appear coherent because it has not yet
encountered implementation constraints, institutional conflict, workload
pressure, or field conditions that would test whether its proposed review
paths are workable across constitutional decisions and decisions within
bounded local governance domains.

The public evidence base for Axodus is currently limited to design
documentation. There are no results with which to assess feasibility,
usability, governance legitimacy, calibrated reliance, contestable challenge
paths, or social impact. The paper reports no field validation, no
institutionally situated validation, no production use, no token-design or
token-distribution claims, no claims about financial operation, no security
guarantees, and no legal or regulatory conclusion.

Human oversight can become nominal when reviewers lack time, authority,
evidence access, or practical power to change the outcome. AI assistance may
be unreliable, difficult to calibrate, or difficult to audit under realistic
workload and time pressure. Weak contestability or poorly designed challenge
paths could make traceability visible without producing answerability,
corrective capacity, effective remedy, or meaningful institutional
accountability.

Participation-readiness mechanisms present particular risks of exclusion,
surveillance, curriculum capture, and false confidence. Modular separation
can also hide systemic coupling, while constitutional constraints can be
overly rigid or too ambiguous for effective local governance. The proposed
architecture may be too complex to govern in practice without substantial
scope reduction.

The synthetic scenario is explanatory rather than empirical. Provenance
vocabulary does not establish record integrity or accountability, and the
proposed measures assess reconstructability, cross-layer inspectability,
review conditions, and bounded reliance rather than legitimacy, fairness,
effective contestability, remedy, or full institutional accountability.
Alternative decompositions may prove clearer or more useful than the six-layer
model. Institutionally situated evaluation may also expose conflicts,
incentives, and workload effects that a controlled artifact cannot reproduce.

The paper does not provide financial, security, investment, or legal advice.
Its terminology may not map cleanly onto every jurisdiction or research
community.

\section{Research Agenda}

Near-term work should prioritize the derivation and interface artifacts,
synthetic scenarios involving bounded local governance domains, and review
tasks suitable for reconstruction, cross-layer coherence, challenge-condition,
and bounded-reliance studies.
Additional literature hardening around governance quality, reviewer
independence, empirical institutional challenge paths, and participation-
readiness validity remains a later requirement before any broader external
review is considered.

\section{Conclusion}

Axodus is proposed as a prototype-stage conceptual architecture for examining
how human-AI-assisted governance might become more traceable,
reconstructable, reviewable, and supportive of accountability and
contestability in federated modular digital institutions. Its contribution is
not a product catalog, implementation-status claim, or assertion of a uniquely
correct decomposition, but an inspectable six-layer arrangement linking
constitutional constraints, bounded local governance domains, functional
service modules, evidence and provenance, advisory AI assistance, and review
paths.

The next research step is to instantiate the conceptual representation in a
controlled experimental environment and test whether independent reviewers can
reconstruct decisions, identify cross-layer inconsistencies, inspect challenge
conditions, and evaluate bounded AI reliance. Those results would test the
architecture's proposed support conditions; they would not independently
establish legitimacy, effective remedy, or institutional accountability.

\renewcommand{\bibpreamble}{%
\noindent References are maintained in \texttt{references.bib}. Verification
records for all cited entries are maintained in the repository's bibliography
control records; final metadata rechecking is required before any submission,
public release, or external circulation.\par\medskip}
\bibliographystyle{plainnat}
\bibliography{references}

\end{document}
